\section{Related works}\label{sec:relworks}

Intrusion detection systems (IDSs) have been extensively studied during last decades, and several surveys on the achievements in this area have been produced \cite{liao2013intrusion,borkar2017survey,hindy2018taxonomy}. 
Some of the surveys available in the state-of-the-art aim their attention at a specific application domain: e.g., some authors review the IDS solutions for Internet of Things \cite{sherasiya2016survey,chaabouni2019network}, Can \etal~ focus on IDSs in wireless sensor networks \cite{can2015survey},  Sultana \etal~ analyse the research works in SDN-based networks \cite{sultana2019survey}. 

IDSs are usually classified into two categories: (i) host-based intrusion detection systems (HIDSs) which monitor and analyse the devices from their internal behavior \cite{liu2018host} and (ii) network-based intrusion detection systems (NIDSs) which monitor and analyse the network traffic to detect any malicious behavior \cite{chaabouni2019network}. In this work, we focus on HIDSs.

In Unix-like operating systems, system calls are the primary source of data used as input for detecting malicious behaviors in HIDSs \cite{liu2018host}. Monitoring the system calls to detect unusual behaviors in computer systems has been done since the nineties \cite{hofmeyr1998intrusion}.

Considering the large amount of calls that are usually generated in a system, multiple algorithms were implemented to extract information carried by the list of system calls. Some common methods are: the n-gram formatting \cite{khreich2017anomaly}, sliding window algorithm \cite{eskin2001modeling}, bag-of-words model \cite{kang2005learning}, and term frequency-inverse document frequency (TF-IDF) method \cite{chen2005application}. All these strategies get a sequence of  system calls as input and output a compressed representation of it.

One of the representation methods that has gained a lot of interest during last decade is the ``Bag of System Calls'' (BoSC) which was introduced for the first time by Kang \etal~ in 2005 \cite{kang2005learning}. Since then, BoSC was widely adopted as an effective and efficient way to represent list of system calls in intrusion and anomaly detection \cite{chen2005application}.



With the recent advancements of cloud systems which enable the execution of virtual applications over rented cloud resources, intrusion detection strategies based on the analysis of system calls at the hosting machine level has gained a lot of interest \cite{ates2018taxonomist,karn2020cryptomining}. 
Thus, a cloud service provider can build a scalable and efficient mechanism to monitor in parallel multiple containers running on the same machine with a single instance of intrusion detection. However, results provided so far show that only few solutions obtain acceptable detection rates.
Ates \etal{} \cite{ates2018taxonomist} developed the ``Taxonomist,'' a technique based on traditional classification models (e.g., random forest, decision trees, and SVM) which is able to distinguish the behavior of different application running on virtualized systems. The proposed mechanism receives as input the information coming from the resource utilization of each application. The proposed classification model is able to recognize unknown applications with high accuracy; however, they did not show what would happen if a malicious code is executed in parallel to the benchmarked applications.
Karn \etal{} \cite{karn2020cryptomining} proposed a supervised approach based on ensemble learning and decision trees for extracting cryptomining signature for explainability. They obtain high accuracy in detecting the cryptoming software behavior, however the solution is based on a supervised model with a focus on cryptomining and it may fail to detect new and unknown malicious behavior.

Finally, a framework for proactive intrusion detection was introduced by Gupta \etal{} \cite{gupta2017supervised}. The proposed model is based on a two stage LSTM-based algorithm aiming at analyzing how the resource utilization varies over the time and classifying normal and abnormal behavior in cloud containers. However, this solution yet obtains high false positive rates (about 10\%), which is undesirable in systems that are monitoring an high number of processes.

Unsupervised approaches have shown interesting results in detecting unknown malicious behavior in several domains \cite{casas2012unsupervised,zhang2006anomaly,zoppi2021unsupervised}. Cui and Umphress \cite{cui2020towards} recently proposed an unsupervised model for detecting attacks in containerized applications. Their solution is based on LSTM autoencoders to detect malicious behaviors. However, their results show high False positive and True positive rates which suggest that much more effort should be put in place to identify unsupervised solutions that would work in cloud environments.
